\section*{Capítulo \ref{ch:b2:two-wave-interference} --- Interferência de duas ondas}

\begin{solution}{exo:b2:two-wave-interference:1}
$i = \lambda D/a = \qty{2.4}{mm}$; $x/i = 4.2$; décima franja escura em $9.5i = \qty{22.5}{mm}$;
$a = \qty{2}{mm}$: \qty{0.47}{mm}; na água $\lambda/n$: \qty{1.8}{mm}.
\end{solution}

\begin{solution}{exo:b2:two-wave-interference:2}
$1.25I_0 \pm I_0$: $V = 0.8$. $1.01I_0 \pm 0.2I_0$: $V = 0.2$ --- o termo de interferência
vai com a razão das \emph{amplitudes}, $1/10$, e não com a das intensidades. Filtro:
$V = 2\sqrt{0.5}/1.5 = 0.94$.
\end{solution}

\begin{solution}{exo:b2:two-wave-interference:3}
$2ne = \qty{665}{nm}$: claro em $665/(p + \tfrac12)$: \qty{443}{nm} (azul-violeta);
escuro em $665/p$: \qty{665}{nm} (vermelho): uma película azul. $e = \qty{1}{\micro m}$: $2ne = \qty{2660}{nm}$,
claro em $760$, $591$, $484$, \qty{409}{nm} --- quatro cores ao mesmo tempo, um
branco rosado e lavado.
\end{solution}

\begin{solution}{exo:b2:two-wave-interference:4}
$r_p = \sqrt{p\lambda R}$: $1.09$, $1.53$, \qty{1.88}{mm}; $p = r^2/\lambda R = 85$ anéis dentro de
\qty{1}{cm}; o centro é escuro: espessura nula e uma inversão de sinal. Água:
raios divididos por $\sqrt n$, $98$ anéis.
\end{solution}

\begin{solution}{exo:b2:two-wave-interference:5}
(a) $i = 0.8$ e \qty{1.4}{mm}; segunda vermelha \qty{2.8}{mm}, terceira violeta \qty{2.4}{mm}.
(b) A ordem $p$ do vermelho em $1.4p$ mm encontra a ordem $p + 1$ do violeta em $0.8(p + 1)$: a partir de
$p = 2$ as ordens se superpõem. (c) $\delta = ax/D = \qty{2.5}{\micro m}$: faltam $\lambda = \delta/(p +
\tfrac12)$: $714$, $556$, \qty{455}{nm}. (d) Centro branco; depois franjas com a
borda interna azul e a externa vermelha; na terceira ordem as cores
se misturam em branco.
\end{solution}

\begin{solution}{exo:b2:two-wave-interference:6}
(a) $b < \lambda L/4a = \qty{59}{\micro m}$. (b) $\pi ab/\lambda L = 6.7$: contraste $|\sin6.7|/6.7
= 0.06$. (c) Uma fonte de um centímetro de largura cobre centenas de interfranjas de
deslocamento: nenhuma franja. (d) $a < \lambda/\theta = \qty{61}{\micro m}$.
\end{solution}

\begin{solution}{exo:b2:two-wave-interference:7}
(a) $a = 2h = \qty{0.4}{mm}$: $i = \lambda D/a = \qty{1.25}{mm}$. (b) Em $x = 0$ os caminhos
são iguais e ainda assim a franja é escura: a reflexão acrescenta $\pi$ (inversão
de sinal na incidência rasante sobre o meio mais denso). (c) $\delta = ax/D <
\qty{50}{\micro m}$: $100$ franjas. (d) O mar é o espelho e a antena no alto do
penhasco é o anteparo: quando a fonte se eleva, as franjas desfilam e
dão sua altura angular.
\end{solution}

\begin{solution}{exo:b2:two-wave-interference:8}
(a) $2d = 42\lambda$: $d = \qty{12.4}{\micro m}$. (b) $100/42 = \qty{2.4}{mm}$. (c) $2nd = p\lambda$:
$56$ franjas. (d) As franjas caminham para a ponta afastada e se espaçam à medida que $d$
diminui; um quarto de franja vale $\lambda/8 = \qty{74}{nm}$ sobre $d$.
\end{solution}

\begin{solution}{exo:b2:two-wave-interference:9}
(a) Só a primeira (ar $\to$ óleo, mais denso); óleo $\to$ água vai para um índice
menor: $\delta = 2ne + \lambda/2$, claro em $2ne = (p + \tfrac12)\lambda$. (b) $2ne =
\qty{1160}{nm}$: reforçados $773$ (borda do visível), \qty{464}{nm} (azul);
cancelados $1160/p$: \qty{580}{nm} (amarelo): azul. (c) $2ne = \lambda/2$: $e = \qty{72}{nm}$.
(d) $\delta = 2ne\cos r$ diminui com o ângulo: os comprimentos de onda reforçados
se deslocam para o azul. (e) As duas reflexões passam então a inverter o sinal, sem
$\lambda/2$ extra: claro em $2ne = p\lambda$, escuro em $(p + \tfrac12)\lambda$ --- a
camada antirreflexo de um quarto de onda.
\end{solution}

\begin{solution}{exo:b2:two-wave-interference:10}
(a) $(n - 1)e = \qty{10}{\micro m}$, $20$ franjas, na direção da fenda coberta (onde
o caminho geométrico deve encurtar \qty{10}{\micro m} para restaurar $\delta = 0$):
$x = (n - 1)eD/a = \qty{2}{cm}$. (b) A franja branca fica ali; a
dispersão do vidro ($\Delta n \approx 0.01$ no visível, ou seja, \qty{0.2}{\micro m} de caminho)
está abaixo de $\ell_c$: ainda visível, ligeiramente colorida. (c) Meça o
deslocamento da franja branca: $e = (\text{deslocamento})\,a/(n - 1)D$. (d) $e(n - 1)\theta^2
(1 + 1/n)/2 = \qty{63}{nm}$ a \qty{5}{\degree}: um oitavo de franja.
\end{solution}

\begin{solution}{exo:b2:two-wave-interference:11}
(a) $\theta = \qty{2.3e-7}{rad}$: $a = 1.22\lambda/\theta = \qty{3.1}{m}$. (b) Seu limite $1.22\lambda/d =
\qty{0.057}{''}$ está acima do diâmetro da estrela, e a atmosfera borra
até $1''$. (c) $1.22\lambda/100 = \qty{7e-9}{rad} = \qty{1.4}{mas}$. (d) As franjas exigem
$\delta < \ell_c = \lambda^2/\Delta\lambda = \qty{32}{\micro m}$.
\end{solution}

\begin{solution}{exo:b2:two-wave-interference:12}
(a) $\delta = 2ne\cos r + \lambda/2$; $\cos r = \sqrt{1 - \sin^2i/n^2} \approx 1 - i^2/2n^2$.
(b) $2ne/\lambda = 10186.8$: o centro, com $p + \tfrac12$ entre $10186.5$ e
$10187.5$, não é nem claro nem escuro. (c) $i_p^2 = 2n^2[1 - (p + \tfrac12)\lambda/
2ne]$ para $p = 10186$, $10185$, $10184$: $i = 0.011$, $0.024$, \qty{0.032}{rad}, $r = fi
= 5.4$, $11.8$, \qty{15.8}{mm}. (d) Só o ângulo $i$ fixa a ordem: cada
ponto da fonte alimenta o mesmo anel sob o mesmo ângulo, de modo que os anéis, vistos
no infinito, sobrevivem a uma fonte larga.
\end{solution}

\begin{solution}{pb:b2:two-wave-interference:1}
\textbf{1.} Cada espelho forma de $S$, por simetria, uma imagem à distância $d$ da
linha de contato; as duas imagens ficam sobre o círculo de raio $d$, separadas
pelo ângulo $2\alpha$: $a = 2\alpha d = \qty{2}{mm}$.

\textbf{2.} $S_1$ e $S_2$ são duas imagens da mesma vibração: seus
saltos de fase são os mesmos.

\textbf{3.} $i = \lambda(d + D)/a = 589 \times 10^{-9} \times 1.7/2 \times 10^{-3} = \qty{0.50}{mm}$: dez
franjas ao longo de \qty{5}{mm}.

\textbf{4.} $\delta = ax/(d + D) = \qty{2.9}{\micro m}$ na borda, $p = 5$: muito abaixo de
$\ell_c/\lambda = 34\,000$.

\textbf{5.} Deslocamento das franjas $bD'/d = 20 \times 10^{-6} \times 1.7/0.2 = \qty{0.17}{mm}$, um terço de
$i$: contraste $|\operatorname{sinc}(\pi ab/\lambda d)| = \operatorname{sinc}(1.07) = 0.82$.

\textbf{6.} $\pi ab/\lambda d = 5.3$: $|\sin5.3|/5.3 = 0.16$.

\textbf{7.} Amplitudes $\sqrt{0.9}$, $\sqrt{0.8}$: $V = 2\sqrt{0.72}/1.7 = 0.998$.

\textbf{8.} $(n - 1)e = \qty{5}{\micro m}$: $8.5$ franjas; acompanhadas sem dificuldade com
luz de sódio; em luz branca a franja branca se desloca de $8.5i = \qty{4.2}{mm}$
e continua ali (a dispersão do vidro custa uma fração de
micrômetro).

\textbf{9.} Uma franja central branca, duas ou três franjas coloridas de
cada lado, depois um branco uniforme.

\textbf{10.} Ar $\to$ óleo inverte, óleo $\to$ água (índice menor) não:
$\delta = 2ne + \lambda/2$; claro para $2ne = (p + \tfrac12)\lambda$.

\textbf{11.} $e = \lambda/4n$: vermelho \qty{112}{nm}, verde \qty{95}{nm}, violeta \qty{72}{nm}.

\textbf{12.} $2ne = \qty{1740}{nm}$: reforçados $1740/(p + \tfrac12)$: $696$ (vermelho), $497$
(azul-esverdeado), \qty{387}{nm}; cancelados $1740/p$: $580$ (amarelo), \qty{435}{nm}:
uma mistura arroxeada de vermelho e azul-esverdeado.

\textbf{13.} A espessura varia ao longo da estrada; cada cor é clara
onde $2ne$ tem o valor certo, de modo que as faixas são contornos de igual
espessura.

\textbf{14.} As ordens diminuem: as cores percorrem a sequência
rumo à extremidade fina; em $e \to 0$ a única inversão de sinal deixa
$\delta = \lambda/2$ para todas as cores: escuro --- como numa película de sabão no ar.

\textbf{15.} Alguns nanômetros: $2ne \ll \lambda$ para todas as cores, as duas
reflexões estão em oposição de fase e se cancelam: preto.

\textbf{16.} $\sin r = \sin60^\circ/1.45$, $r = \qty{37}{\degree}$, $\cos r = 0.80$: $2ne\cos r
= \qty{1390}{nm}$, claro em $928$, $557$, \qty{398}{nm}: o vermelho da questão 12
virou verde-amarelado --- rumo ao azul.

\textbf{17.} Quando $2ne$ ultrapassa $\ell_c \approx \qty{1}{\micro m}$, ou seja, além de uns
\qty{0.4}{\micro m} de óleo, as ordens se superpõem e as cores se apagam num cinza.

\textbf{18.} Os dois raios refletidos pela película, para qualquer direção
da fonte, se reencontram sobre a película: as franjas ficam localizadas ali, e
uma fonte larga e multidirecional apenas soma padrões idênticos.

\textbf{19.} $\ell_c/\lambda$: $34\,000$ com o sódio, $1.7 \times 10^7$ com o laser.

\textbf{20.} $b < \lambda d/4a$: \qty{15}{\micro m} para os espelhos, \qty{59}{\micro m} para
as fendas.

\textbf{21.} Desenhe dois pontos-fonte: para cada um, os dois raios refletidos
se recombinam no mesmo ponto da película, com o mesmo $\delta = 2ne\cos r$
sob incidência quase normal; o padrão da película não se move, o de Young
se move.

\textbf{22.} $\lambda/\Delta\lambda = 589/0.02 \approx 30\,000$.

\textbf{23.} $550/30 \approx 18$ franjas.

\textbf{24.} A luz coerente espalhada pelos grãos do anteparo interfere
com fases aleatórias no olho: o speckle.

\textbf{25.} Espelhos: \omterm{prop:b2:two-wave-interference:young}{divisão da frente de onda}, franjas em toda parte
(não localizadas), número limitado por $\ell_c/\lambda$, uma fonte larga as
apaga. Película: \omterm{prop:b2:two-wave-interference:film}{divisão da amplitude}, franjas sobre a película, cores
limitadas por $\ell_c \sim \qty{1}{\micro m}$, uma fonte larga é inofensiva.
\end{solution}
